% Proposed native author continuation. Integrator chooses the module number.
\section{Clearance-selected endpoint kernels and critical boundary tightness}
\label{sec:ce-kernel}

The endpoint layers left by Theorem~\ref{thm:lc-deep-payment} cannot
be discarded by their collision-area mass.  We establish a different
physical estimate: a clearance source at any flight has a small
complete endpoint-momentum kernel after propagation along its longer
side.  This pays the first outgoing flight as well as the later
flights.  It subsequently gives a uniform tightness estimate for
positive-clearance critical sectors.  The norm of that conclusion
and its distinction from the original roof-height norm are part of
the statement.

For $1\le j\le m-2$, the selected endpoint estimate below
already follows from Theorem~\ref{thm:lc-local-half} in the two
physical time directions. The additional first-image estimate
retains a witness selected before the first true collision and
thereby completes $j=0$ and $j=m-1$. The longer-side argument is
also present in the retained zero-clearance fan theorem
\ref{pc:fan}; that theorem already pays its entire zero-level
stratum at the original collision normalization. The new critical
conclusion here concerns a positive clearance width, with the
different norm explicitly displayed below.

All constants and curves are those of Section~\ref{sec:lc-length}.
In particular the original radius range, the actual first-hit map,
the two impact signs and every unvisited disk are retained.  Put
\begin{align}
 c_H&=\sin(\widehat k_0^{-2}),\qquad
 C_H=\max\{c_H^{-1},(\pi/2)(1+\widehat k_0^{-1})^2\},\notag\\
 D_{\rm first}
   &=C_H\,\frac{1+\tau_+/R_-}{t_*/R_+}\sqrt{1+U^2},
                                                        \label{eq:ce-first-ratio}\\
 S_{\rm first}&=K_{\rm c}A_{\rm hor},\qquad
 M_{\rm first}=D_{\rm first}J_0L_{\rm c}/R_-,\notag\\
 C_{\rm first}
   &=\frac{J_f}{k_-t_*}
       \{\Gamma S_{\rm first}/\theta_{\rm r}
                                      +C_\ell M_{\rm first}\},\notag\\
 C_{\rm end}
   &=\max\{K_{\rm loc},C_{\rm first},J_0J_fE_{\rm lc}\}.
                                                        \label{eq:ce-constants}
\end{align}
These constants are finite and independent of $m,R,j,\delta,a,v$.
The enclosing physical flight bound $\tau_+$, rather than a sampled
maximum, occurs in \eqref{eq:ce-first-ratio}.

\subsection{Weighted first images of a selected horizontal source}

\begin{lemma}[The first selected flight retains its length]
\label{lem:ce-first}
For a fixed horizontal source $F_a$, let its initial outgoing
flight satisfy the original clearance event $E_{\delta,R}$,
where $0<\delta\le\delta_*$.  A positive cover of its first
physical image by homogeneous short labelled curves $C$ carries
the exact image densities $\rho_C$ and satisfies
\begin{equation}\label{eq:ce-first-budgets}
 \sum_C\|\rho_C\|_\infty\le S_{\rm first},\qquad
 \sum_C\|\rho_C\|_\infty|C|\le M_{\rm first}\delta .
\end{equation}
The cover may include further initial states, but each of them
follows its actual first collision.  There is no reflection at a
clearance foot.
\end{lemma}
\begin{proof}
Divide $F_a$ into the $J_0$ short preliminary arcs of
\eqref{eq:lc-source-counts}.  On each arc,
Lemma~\ref{lem:lc-cover} covers its selected points by at most
$K_{\rm c}$ disjoint initial intervals.  The combined
$d\alpha_0$ mass of all these intervals is at most
$J_0L_{\rm c}\delta/R_-$.

For each initial interval take its actual first-hit pieces and
their first homogeneous images.  Cut images into the same short
lengths as in the proof of \eqref{eq:lc-source-budgets}.  An initial
subinterval cannot increase the number of first-hit pieces on
its preliminary arc, and on each first-hit piece the first-image
momentum angle is strictly monotone.  Consequently each target
homogeneity strip is still met at most once.  Applying the proof
of the $A_{\rm hor}$ budget to at most $K_{\rm c}$ intervals per
preliminary arc gives the first inequality in
\eqref{eq:ce-first-budgets}.  Source labels distinguish any
coincident target coordinate images.

For the second inequality use the exact density of
Lemma~\ref{lem:mc-first},
\[
 \rho_C=\frac{d}{R A\sqrt{1+V_1^2}},\qquad
 A=\sqrt{1-a^2}+\tau/R,\qquad
 d=\cos\varphi_1 .
\]
On a central homogeneous image, the ratio of the largest and
smallest $d$ is at most $c_H^{-1}$.  On a signed strip of index
$k\ge\widehat k_0$ it is at most
\[
 \frac{\sin(k^{-2})}{\sin((k+1)^{-2})}
       \le\frac{\pi}{2}\left(\frac{k+1}{k}\right)^2
       \le\frac{\pi}{2}(1+\widehat k_0^{-1})^2 .
\]
Also $t_*/R_+\le A\le1+\tau_+/R_-$ and
$1\le\sqrt{1+V_1^2}\le\sqrt{1+U^2}$.  Thus every such
image and each of its subcurves satisfies
\[
 \sup_C\rho_C\le D_{\rm first}\inf_C\rho_C,\qquad
 \|\rho_C\|_\infty|C|
                \le D_{\rm first}\int_C\rho_C\,d\ell_C .
\]
The formulas hold on open regular pieces; taking the endpoint
limits within one homogeneity strip gives the same inequalities.
Their exact integrals sum to the $d\alpha_0$ mass of the initial
cover.  The latter intervals are disjoint after merging within
each preliminary arc.  The first-hit and homogeneous cuts have
zero source measure, so summing proves the second inequality.
The countable strip sum is legitimate by nonnegativity.  No
bounded distortion across different grazing strips is asserted.
\end{proof}

\subsection{The complete selected kernel}

Write $\mathcal R_{m,R}(a,v)$ for the unique regular physical
roots in \eqref{eq:mc-roots}.  At such a root put
$u_x=|\partial_{\alpha_0}p_m(x)|^{-1}$ and define
\begin{equation}\label{eq:ce-root-kernel}
 D_{m,R;j,\delta}(a,v)
   =\sum_{\substack{x\in\mathcal R_{m,R}(a,v)\\
                         N_{\delta,R}(T_R^jx)\ge1}}u_x,
 \qquad 0\le j<m .
\end{equation}
The root convention counts a physical state once, including
states on artificial homogeneity or coordinate seams.

\begin{theorem}[Clearance-selected endpoint-momentum kernel]
\label{thm:ce-kernel}
For all $m\ge1$, $R\in[R_-,R_+]$, $a,v\in(-1,1)$,
$0\le j<m$ and $0<\delta\le\delta_*$,
\begin{equation}\label{eq:ce-kernel}
 D_{m,R;j,\delta}(a,v)
       \le C_{\rm end}
          \{\delta+\theta_{\rm r}^{\max(j,m-1-j)}\}.
\end{equation}
The corresponding selected pushforward of $d\alpha_0$ on $F_a$
has this density bound.  Therefore, with the original conditioning
and any original Borel selector $0\le H\le1$,
\begin{equation}\label{eq:ce-joint}
 (p_0,p_m)_\#
   \bigl(H\one_{E_{j,\delta,R}}\,d\nu_R^*\bigr)
 \le \frac{C_{\rm end}}{4\pi c_*}
       \{\delta+\theta_{\rm r}^{\max(j,m-1-j)}\}\,da\,dv .
\end{equation}
All original roof, return, displacement, collision and observation
tests may be retained by positive restriction.  These inequalities
are a coordinate disintegration of the real orbit law and do not
make the momentum process Markov.
\end{theorem}
\begin{proof}
First prove the forward bound with exponent $m-1-j$.  If
$1\le j<m$, apply Theorem~\ref{thm:lc-local-half} to the
horizontal source with
\[
 h=j-1,\qquad n=m-j-1,\qquad l=h+n+2=m .
\]
Its inverse-derivative supremum sum is at most
$K_{\rm loc}(\delta+\theta_{\rm r}^{m-j-1})$.  On each
source interval the terminal momentum is strictly monotone,
so it contributes at most one root at a specified terminal
value.  The same sum bounds its selected density almost
everywhere by one-dimensional coarea.

For $j=0$ and $m\ge2$, use Lemma~\ref{lem:ce-first}.
Propagate each first-image curve through exactly $m-2$ further
actual collisions and then take the final collision without an
infinite terminal homogeneity cut.  The renewal inequality
\eqref{eq:lc-renewal}, weighted by $\|\rho_C\|_\infty$,
gives a sum of incoming density suprema at most
\[
 \Gamma\theta_{\rm r}^{m-2}
          \sum_C\|\rho_C\|_\infty
            +C_\ell\sum_C\|\rho_C\|_\infty|C|
 \le \Gamma S_{\rm first}\theta_{\rm r}^{m-2}
                          +C_\ell M_{\rm first}\delta .
\]
The exact last-collision derivative cancellation
\eqref{eq:mc-last-cancellation} and at most $J_f$ actual
last-hit pieces multiply this bound by $J_f/(k_-t_*)$.
Since $\theta_{\rm r}^{m-2}
=\theta_{\rm r}^{-1}\theta_{\rm r}^{m-1}$, this is at most
$C_{\rm first}(\delta+\theta_{\rm r}^{m-1})$.
The generation count is $1+(m-2)+1=m$.
For $m=1$, the unselected root sum is at most
$J_0J_fE_{\rm lc}$ by the first-hit count and its exact
derivative lower bound.  This proves the asserted forward
density bound also in that case.

We justify the bound at every regular root value, including
partition seams.  Take any finite selected set of distinct
regular roots at $(a,v)$.  A genuine width-$\delta$ witness
has strict interior-foot and strict clearance inequalities;
it persists in a sufficiently small regular neighborhood
of each root.  The inverse branches and positive weights
are continuous.  Their sum is bounded almost everywhere at
nearby terminal values by the density estimate just proved,
so it obeys the same bound at $v$.  Exhausting finite root
sets gives the complete series.  A weak width-$\delta$ version,
when needed, follows by using $\delta+\eta\le\delta_*$
and letting $\eta\downarrow0$; at $\delta=\delta_*$ one can
instead retain the strict convention of the theorem.

Now reverse each actual finite record.  The collision
reverser is $(\alpha,p)\mapsto(\alpha,-p)$, the selected
flight has index $m-1-j$, and the two endpoint momenta
become $(-v,-a)$.  A flight of length $\tau$ is the same
unoriented physical segment after reversal.  For a designated
unvisited disk its closest-foot coordinate changes from
$t_{\rm foot}$ to $\tau-t_{\rm foot}$, so the strict
interior-foot condition and its unsigned clearance are
preserved.  The two incident disks are exchanged and the
same third disks remain unvisited on this flight.  The
original terminal disk has one definite lattice center;
subtracting that center re-centers the reversed starting
disk in the original fundamental cell uniquely.  This
bijection introduces no lattice or word multiplicity.
Every original observation is transported on this
same reversed orbit, rather than assumed invariant.
Area preservation of the physical
collision map gives equality of the original and reversed
inverse root weights.  Explicitly, on a regular branch write
$DT_R^m$ in the coordinates $(\alpha,p)$.  Its determinant
is one; the lower-left entry of its inverse has the opposite
sign and the same magnitude as the original lower-left
entry.  The two momentum sign changes preserve that
magnitude.  Thus the reversed forward bound is
$C_{\rm end}(\delta+\theta_{\rm r}^{j})$.
Taking the smaller of the two bounds, using
$0<\theta_{\rm r}<1$, proves \eqref{eq:ce-kernel}.
The two sides were evaluated on the same trajectory;
no product of their probabilities was used.

Finally integrate the fibre inequality in $a$ and restore
the original positive selectors.  The collision-coordinate
measure is $d\alpha_0\,da/(4\pi)$ and the original section
conditioning contributes exactly $c_*^{-1}$.  This proves
\eqref{eq:ce-joint}; it is not a rescaling by the selected
source mass.
\end{proof}

\subsection{Uniform tightness of positive-clearance critical sectors}

Consider the physical one-copy critical sectors of
Section~\ref{sec:mc-critical}, using the finite analytic
side conditions justified in Lemma~\ref{pc:representatives}.
All visited incidences at the limiting germ are positive;
unvisited tangencies, finite original section cuts and corners
are allowed.  Keep the original fractions $\beta_y(x)$,
$\beta_z(x)$ and actual critical coefficient $\mathfrak J(x)$.
At a marked limiting flight count a witness of width $\delta$
when its strict auxiliary flight has an interior closest foot
and clearance in $[0,\delta]$ from an unvisited disk.  This
definition is used only for its actual physical-side germ.
It does not insert a collision at clearance zero.

\begin{theorem}[Positive-clearance critical boundary tightness]
\label{thm:ce-critical}
For $0<\delta<\delta_*$, the complete one-copy family
$\mathcal F_{m,j,\delta}(R)$ of those critical germs with
at least one specified-flight witness obeys
\begin{equation}\label{eq:ce-critical}
 \sum_{x\in\mathcal F_{m,j,\delta}(R)}
          \mathfrak J(x)\beta_y(x)
 \le C_{\rm crit}
       \{\delta+\theta_{\rm r}^{\max(j,m-1-j)}\},
 \qquad C_{\rm crit}=\frac{U_2C_{\rm end}}{2c_*}.
\end{equation}
A germ with several witnesses is counted once in this sum.
Original exact-record and critical-roof-value restrictions,
and further weights in $[0,1]$, may be retained.

For the original clearance guard
$\delta_{j,B}=2\varepsilon(B)q_0^{\min(j,m-1-j)/4}$
with $2\varepsilon(B)<\delta_*$, the leading critical
coefficient of the union of all its clearance witnesses is
at most
\begin{equation}\label{eq:ce-guard-critical}
 C_{\rm crit}\left\{
   \frac{4\varepsilon(B)}{1-q_0^{1/4}}
                 +m\theta_{\rm r}^{(m-1)/2}\right\}.
\end{equation}
Thus its unscaled complete critical coefficient is uniformly
tight as the original guard is removed, after the collision
limsup and the original radius and exact-label suprema.
\end{theorem}
\begin{proof}
Fix $R,m,j$ and a finite set of distinct one-copy critical
germs.  Use the common small endpoint-momentum disk and
their inverse branches from the proof of
Theorem~\ref{thm:mc-critical-sectors}.  On an actual
regular side of a germ whose limiting witness has width
at most $\delta$, that witness has width less than
$\delta+\eta$ for every fixed $\eta>0$ on a sufficiently
small neighborhood.  Its strict foot remains interior by
the uniform standoff from the incident contacts.  If the
limiting clearance is zero, the actual unvisited side has
positive clearance there.  Choose
$\eta>0$ with $\delta+\eta\le\delta_*$.

The physical-side inverse images are disjoint after the
original one-copy merging.  Use the unconditioned
coordinate measure $d\alpha_0\,dp_0$ here.  Its
endpoint-momentum pushforward has density given by the
regular fibre root sum \eqref{eq:ce-root-kernel}, by
one-dimensional coarea in $\alpha_0$ followed by
integration in $p_0$.  Theorem~\ref{thm:ce-kernel}
therefore gives almost everywhere in $z=(p_0,p_m)$
\[
 \sum_xu_x(z)\one_{E_x}(z)
     \le C_{\rm end}
        \{\delta+\eta+\theta_{\rm r}^{\max(j,m-1-j)}\}.
\]
The finite analytic physical-side conditions have angular
limits outside finitely many rays, by
Lemma~\ref{pc:representatives}.  Polar Fubini, continuity of
the finite inverse weights and bounded convergence along
good radii tending to zero therefore yield
\[
 \sum_xu_x\beta_z(x)
     \le C_{\rm end}
        \{\delta+\eta+\theta_{\rm r}^{\max(j,m-1-j)}\}.
\]
Let $\eta\downarrow0$, then exhaust finite germ families.
These limits occur at the fixed physical $R,m,j$;
their resulting bound has constants uniform in those
parameters.  Finally the exact angular comparison
\eqref{eq:mc-sector-comparison} and
$\lambda_{\max}(G_x)/R^2\le U_2$ give
\eqref{eq:ce-critical}.  The only section-conditioning
factor is the $c_*^{-1}$ already present in that critical
coarea comparison; none was used in the unconditioned
kernel inequality.  This passage controls actual
one-sided coefficients, rather than assigning zero trace
to a collision-area-null set.

For the original union, positivity bounds its one-copy
coefficient by the sum of the $m$ selected coefficients.
There are at most two indices at each depth, so
\[
 \sum_{j=0}^{m-1}\delta_{j,B}
       \le\frac{4\varepsilon(B)}{1-q_0^{1/4}},\qquad
 \sum_{j=0}^{m-1}\theta_{\rm r}^{\max(j,m-1-j)}
       \le m\theta_{\rm r}^{(m-1)/2}.
\]
Insert these finite-count bounds into
\eqref{eq:ce-critical}.  Complete positive sums precede all
radius, label and roof-value restrictions.  Taking the
collision limsup first, and only then $B\to\infty$ along
the original $\varepsilon(B)\to0$, proves the stated
unscaled tightness.
\end{proof}

\paragraph{The resulting gain and the exact limitation.}
Theorem~\ref{thm:ce-critical} treats positive-width
clearance witnesses at every flight, including both
shallow endpoint layers.  The previous zero-clearance
coefficient payment is retained and is sharper on its
own zero-level stratum.  The new argument retains
selected physical length through the first collision
and then uses the longer genuine propagation direction.
It supplies a uniform small-source estimate in the
endpoint-momentum kernel and in the complete critical
boundary coefficient norm.

It does not bound the original shortage in
$m^2L^\infty(dt)$.  In particular multiplying the first
term of \eqref{eq:ce-guard-critical} by $m^2$ at a fixed
$B$ gives no vanishing bound.  A leading critical
coefficient is also not a collision-uniform estimate of
the nonleading coarea remainder at a common original
roof.  Those remainders, the original same-pool
endpoint-layer shortage, and the joint spectral,
arithmetic and conditioning inputs remain as specified
in the unchanged main theorem and dependency ledger.
No neighborhood radius uniform over the full critical
family has been inferred from finite-germ exhaustion.
